%%%%%%%%%%  scan idx 19  |  folio 6  %%%%%%%%%%
\SGAPrintPage{6}
\sgaStableTarget{sga6.e00.sec.3}\subsection*{3.}
Let us now show how we can still give a meaning to formula (1.2) when we abandon the regularity hypothesis on $X,Y$, replacing it by a hypothesis of a local nature on the morphism $f$, which we shall specify.

\sgaStableTarget{sga6.e00.sec.3.1}\subsubsection*{3.1}
A first difficulty arises because the observation that served as the starting point of the theory [RRR], namely that the group $K(X)$ of classes of sheaves on $X$ is the same whether it is defined in terms of coherent sheaves or in terms of locally free sheaves, is essentially tied to the regularity hypothesis on $X$. When $X$ is a general scheme (say noetherian, for simplicity), we must distinguish these two groups, which we shall denote respectively by $K_\bullet(X)$ and $K^\bullet(X)$, the first being covariant in $X$ for proper morphisms, and the second contravariant in $X$ for arbitrary morphisms. We note that $K^\bullet(X)$ has a natural ring structure (and even an augmented $\lambda$-ring structure), whereas $K_\bullet(X)$ in general no longer has a ring structure, but only a module structure over $K^\bullet(X)$. The homomorphism $x\to x\,\cl_\bullet(\underline O_X)$ (where $\cl_\bullet(\underline F)$ denotes the class of $\underline F$ in $K_\bullet(X)$) from $K^\bullet(X)$ to $K_\bullet(X)$ is in general no longer an isomorphism. \uline{For our formulation of the generalized Riemann--Roch theorem, we shall work exclusively with $K^\bullet(X)$, $K^\bullet(Y)$, to the exclusion of $K_\bullet(X)$, $K_\bullet(Y)$}. As we noted in 2.3., these augmented $\lambda$-rings, equipped with their natural $\lambda$-filtration, give rise to graded rings $\Gr^\bullet(X)$, $\Gr^\bullet(Y)$, and ``Chern class'' homomorphisms
\[
c_i:K^\bullet(X)\longrightarrow \Gr^\bullet(X),
\]
and similarly for $Y$. The difficulty, however, arises in defining a direct-image homomorphism
\sgaStableTarget{sga6.e00.eq.3-1}\begin{equation}\tag{3.1}
f_*:K^\bullet(X)\longrightarrow K^\bullet(Y),
\end{equation}
which can no longer be defined by the naive formula
\[
f_*(\cl^\bullet(\underline F))=\sum_i(-1)^i\cl^\bullet\bigl(R^if_*(\underline F)\bigr),
\]
because, even when $\underline F$ is locally free, the Modules $R^if_*(\underline F)$ are in general no longer locally free, nor even of finite tor-dimension, and the preceding formula loses its meaning. A solution to this difficulty is provided by the point of view of \underline{derived categories}. Indeed, although the $R^if_*(\underline F)$ are in general ``bad'' sheaves, say of infinite tor-dimension,

%%%%%%%%%%  scan idx 20  |  folio 7  %%%%%%%%%%
\SGAPrintPage{7}
to which we cannot associate a class in $K^\bullet(Y)$, the complex $Rf_*(\underline F)$ on $Y$, the direct image of $\underline F$ in the sense of derived categories, whose cohomology sheaves are the $R^if_*(\underline F)$, nevertheless tends to be excellent. More precisely, suppose that $\underline F$ has finite tor-dimension over $\underline O_Y$, which is the case when $\underline O_X$ has finite tor-dimension over $\underline O_Y$ and, in addition, $F$ is locally free. We then show (Exp. III) that the complex $Rf_*(F)$ is ``\underline{perfect}'', by which we mean here: it has coherent cohomology sheaves (which we already knew in any case) \underline{and} finite global tor-dimension; or, equivalently: on every affine open $U$ of $Y$, this complex is isomorphic, in the derived category $D(U)$, to a complex bounded in degree and locally free of finite type in each degree. From the point of view of Homological Algebra, perfect complexes are ``as good'' as locally free Modules, of which they are the natural generalization. They moreover form a ``triangulated subcategory'' of the derived category $D(Y)$, i.e. a subcategory stable under cones (or ``mapping cylinders'') of morphisms. On the other hand, to every triangulated category $C$ we associate a group $K(C)$ of classes, a natural variant of the construction of the group of classes of an abelian category, by taking the free group generated by the objects of $C$, modulo the relations
\[
\cl(K)-\cl(K')-\cl(K'')
\]
for every distinguished triangle
\[
\begin{tikzcd}[column sep=large]
 & K'' \arrow[dl] & \\
K' \arrow[rr] & & K \arrow[ul] .
\end{tikzcd}
\]
What replaces the fundamental observation of [RRR] in the present case is the fact that, when the scheme $Y$ admits an ample invertible Module, every perfect complex $L_\bullet$ on $Y$ is globally (and not only locally) isomorphic in $D(Y)$ to a complex $L'_\bullet$ bounded in degree and locally free of finite type in each degree (Exp. II); it follows readily (Exp. IV) that the canonical homomorphism
\[
K^\bullet(Y)\longrightarrow K(\Parf(Y)),
\]
where $\Parf(Y)$ is the triangulated category of perfect complexes on $Y$, is an isomorphism. Thus to every perfect complex $L_\bullet$ on $Y$ we can associate a class
\[
\cl^\bullet(L_\bullet)\in K^\bullet(Y),
\]

%%%%%%%%%%  scan idx 21  |  folio 8  %%%%%%%%%%
\SGAPrintPage{8}
which is moreover simply
\[
\cl^\bullet(L_\bullet)=\sum_i(-1)^i\cl^\bullet(L'_i).
\]
where $L'_\bullet$ is as above, so that the $L'_i$ are indeed locally free Modules on $Y$ whose classes we can take in $K^\bullet(Y)$.

With this understood, we define a homomorphism (3.1) whenever $f:X\to Y$ is a \underline{proper morphism of finite tor-dimension} of noetherian schemes such that $Y$ admits an ample invertible Module, simply by putting
\sgaStableTarget{sga6.e00.eq.3-2}\begin{equation}\tag{3.2}
f_*(\cl^\bullet(\underline F))=\cl^\bullet(Rf_*(\underline F))
\end{equation}
for every locally free Module $\underline F$ on $X$ (so that $Rf_*(\underline F)$ is indeed a perfect complex on $Y$, and the right-hand side is defined). When $X$ also admits an ample invertible Module, $K^\bullet(X)$ can likewise be interpreted as $K(\Parf(X))$, and the homomorphism (3.1) is precisely the one associated, by functoriality, with the exact functor of triangulated categories induced by
\[
Rf_*:\Parf(X)\longrightarrow \Parf(Y).
\]

\sgaStableTarget{sga6.e00.sec.3.2}\subsubsection*{3.2.}
To give a meaning to (1.2), we must still impose on $f$ a hypothesis allowing a meaning to be given to the homomorphism (2.5)
\sgaStableTarget{sga6.e00.eq.3-3}\begin{equation}\tag{3.3}
f_*:\Gr^\bullet(X)_{\mathbb Q}\longrightarrow \Gr^\bullet(Y)_{\mathbb Q},
\end{equation}
i.e. ensuring that the homomorphism (3.1) just defined is (up to a shift and modulo torsion) compatible with the filtrations of these rings; and finally we must still define a ``virtual relative tangent bundle'' element
\[
T_f\in K^\bullet(X).
\]
For the latter, returning to the definition proposed in 2.1, we see that we should make a hypothesis ensuring that (with the notation of loc. cit.) the conormal Module $\underline N_{X/X'}$ is locally free. The most natural hypothesis for this purpose is that $i$ is a ``regular immersion'', i.e. identifies $X$ with a subscheme of $X'$ defined by an Ideal which locally can be defined by an $\underline O_{X'}$-regular sequence of sections of $\underline O_{X'}$. This hypothesis, of a purely local nature, does not in fact depend on the choice of the factorization of $f$ into an immersion followed by a smooth morphism (Exp. VIII), and, as above, the same holds for the class $T_f$ defined by formula (2.2). When the hypothesis just considered holds, we say that the morphism $f$ is a ``\underline{locally complete intersection}'' \underline{morphism}\SGAAuthorFootnote{raisedstar}{We shall also allow ourselves to omit the word ``locally'' in this terminology, which entails no risk of confusion.}. Such a morphism is moreover locally of finite tor-dimension.

%%%%%%%%%%  scan idx 22  |  folio 9  %%%%%%%%%%
\SGAPrintPage{9}
For a locally complete intersection morphism $f$, we introduce a natural notion of ``virtual relative dimension'' $d(f)$ by the formula
\[
d(f)(x)=\dim_x\bigl(f'^{-1}f(x)\bigr)-\codim_x(X,X').
\]
This is an integer-valued function, locally constant in the point $x\in X$, hence constant when $X$ is connected; it likewise does not depend on the choice of the factorization (2.1). When this virtual relative dimension $d$ is constant, $f$ is proper, and $X$ and $Y$ admit ample invertible Modules, we can show (Exp. VIII), although this is far from a trivial result, that
\sgaStableTarget{sga6.e00.eq.3-4}\begin{equation}\tag{3.4}
f_*\bigl(\operatorname{Filt}^iK^\bullet(X)_{\mathbb Q}\bigr)\subset \operatorname{Filt}^{i-d}K^\bullet(Y)_{\mathbb Q},
\end{equation}
hence a homomorphism (3.3) of degree $-d$.

\sgaStableTarget{sga6.e00.sec.3.3}\subsubsection*{3.3.}
We have thus assembled all the structural elements needed to give a meaning to formula (1.2) \uline{for a proper locally complete intersection morphism of noetherian schemes $X,Y$ both admitting ample invertible Modules}. With a little additional care, this formulation moreover remains meaningful independently of any noetherian hypothesis.

In this general form, this Riemann--Roch formula will be proved in Exp. VIII. The reader will see that the proof given in this Seminar combines elements of the two proofs in [RRR] which were discussed in the Introduction. The need to use the methods of the first of these proofs, involving more $\lambda$-formalism than the second (followed in Borel--Serre), is due to the need for a proof of inclusion (3.4), a question that did not arise in the more restricted setting of [RRR] and Borel--Serre [2], where we could work with the Chow ring or the ring $\Gr^\bullet_{\mathrm{top}}$. From the second proof, we retain the introduction of the blow-up of $Y$ along $X$ when $f$ is an immersion. This technical device will doubtless disappear, and the proof for an immersion will extend to the case where we no longer postulate the existence of an ample invertible Module on $Y$, once the question raised in Exp. XIV no.~1 concerning the introduction of exterior-power operations in the derived category $D^-(Y)$ has been resolved.

\sgaStableTarget{sga6.e00.sec.4}\subsection*{4.}
Finally, we must give some indications of how a meaning can still be given to (1.2) without a hypothesis of the existence

%%%%%%%%%%  scan idx 23  |  folio 10  %%%%%%%%%%
\SGAPrintPage{10}
of ample invertible Modules on $X$ and $Y$.

\sgaStableTarget{sga6.e00.sec.4.1}\subsubsection*{4.1.}
First, we note that \uline{for an arbitrary scheme $X$ (or, more generally, locally ringed topos), it seems beyond doubt that the ``right'' invariant to replace the ring $K(X)$ of [RRR]} is not \uline{the ring $K^\bullet(X)$ considered above, defined in terms of locally free Modules on $X$, but the group $K(\Parf(X))$, considered in 3.1., defined in terms of perfect complexes on $X$, a group that we shall denote here by} $K^\bullet(X)$, and which we shall distinguish from the ``naive'' group of 3.1, denoted by $K^\bullet_{\mathrm{naive}}(X)$. This point of view is forced by the need to define a direct-image homomorphism (3.1) for a proper morphism $f:X\to Y$ of finite tor-dimension between, say, noetherian schemes (the non-noetherian case requires somewhat stronger local finiteness properties of $f$, cf. Exp. III). With the definition of $K^\bullet(X)$ now adopted, this homomorphism $f_*$ is indeed defined trivially by formula (3.2), where $\underline F$ is now an arbitrary perfect complex on $X$, which implies that $Rf_*(\underline F)$ is also perfect.

\sgaStableTarget{sga6.e00.sec.4.2}\subsubsection*{4.2}
This new definition of $K^\bullet(X)$ nevertheless raises a new problem of homological algebra: that of defining a $\lambda$-structure on $K^\bullet(X)$ compatible with its natural ring structure arising from the tensor product in the derived category. (NB When $L_\bullet$ and $L'_\bullet$ are two perfect complexes, their total tensor product $L_\bullet\otimes^{\mathbf L}L'_\bullet$ is also perfect). More precisely, we should define ``exterior-power'' operations in the derived category $D^-(X)$, taking perfect complexes to perfect complexes. This question has not yet been clarified at present, and, as was stated in the Introduction, it constitutes the most serious gap in the Seminar. Once this question has been resolved, and proceeding as in 2.3, we obtain a natural filtration on $K^\bullet(X)$, allowing us to form the associated graded ring $\Gr^\bullet(X)$, which will play the role of the Chow ring, and to define a theory of Chern classes
\[
c_i:K^\bullet(X)\longrightarrow \Gr^\bullet(X).
\]

\sgaStableTarget{sga6.e00.sec.4.3}\subsubsection*{4.3}
To give a meaning to formula (1.2) for a proper locally complete intersection morphism $f:X\to Y$ of noetherian schemes, we must still include in the statement of this formula the nontrivial inclusions (3.4) (where $d$ is the virtual relative dimension) to allow (3.3) to be defined, and finally define the virtual relative tangent class

